\section{A companion equation in ten positive variables}
\label{sec:compactcompanion}

A second construction uses a different tradeoff: ten strictly positive
unknowns, 78 binary arithmetic operations, and exact total degree 454.
It retains the whole four-integer native coefficient tuple of the incoming
\emph{Centered Residues} package. Its matched parent has 76 operations and
degree 464; that package also supplies 77/460. Thus this companion spends
two or one extra operations to reduce degree by ten or six. This is a
specific comparison, not a global simplicity claim. The new argument
recovers the first Pell square only after recovering the main Pell square.

\subsection{The complete polynomial}

The variables are
\[
 J,F,\alpha,i,T,w,s,n,\eta,\zeta\in\mathbb Z_{>0}.
\]
The fixed ports are \(\beta,K_0,M_C,M_F\); their complete literal bytes
and finite reconstruction are supplied in the accompanying compact
component. They obey \(\beta\ge31\) odd, \(K_0>0\), and
\(0<M_C<\beta<M_F<2\beta\). They are constants, not unknowns.
Define the following polynomial abbreviations:
\begin{align*}
 h&=\beta J,&q&=h+1,&C&=q-F-\alpha,\\
 X&=wq,&Y&=q^3s,&k&=\eta+\zeta,\\
 c&=kY+\eta,&A&=Y(X+1)+1,&\Delta&=A^2-1,\\
 H&=4A-5,&\ell&=(XY)(kY),&E&=1+\ell(\ell+k),\\
 R&=[h(q-F)+\alpha](q^2-1)+(M_C+qM_F)J,\\
 e&=k-R-1,& B&=(K_0+w)C-F.
\end{align*}
The new strong-pair subgraph is
\begin{align*}
 D_m&=1+\Delta c^2,&U&=iAc^2,&S&=ED_m,\\
 a&=\Delta U^2S,&F_2&=1+a,&Q&=\Delta a,&G&=SF_2.
\end{align*}
Complete the equation with
\begin{align*}
 b&=cT-R,&V&=F_2b-c,&L&=bV,\\
 K&=X(3c-2X)+2,&M&=hXYH,\\
 Z&=hXY\{H(n-LB)+K\}-H\{XYB+he\},\\
 P&=Q\{Z^2-M^2GV^2\}+(M^2G-Z^2).
\end{align*}
The single equation is \(P=0\). Every abbreviation is polynomial. In
particular \(F_2\) being a square will be a conclusion; its name is not
a square-root instruction. The graph has exactly ten variable leaves
and four fixed coefficient leaves.

\begin{theorem}
Composed with the retained whole-tuple native/refutation calibration,
ordinary mathematics formalizable in ZFC proves
\[
 \bigl(\forall\mathbf x\in\mathbb Z_{>0}^{10},\ P(\mathbf x)\ne0\bigr)
 \quad\Longleftrightarrow\quad\operatorname{Con}(\mathrm{ZFC}).
\]
The polynomial has exact degree 454 and a 78-operation evaluator.
\end{theorem}

\subsection{Delayed recovery of the first square}

For proof notation set \(\delta=Q-1\), \(C_y=QV^2-1\), and
\(z=Z/M\). Positive inputs give \(F_2>2c\), \(Q>F_2>1\), \(M>0\),
and \(|V|>1\), regardless of the initial sign of \(b\). Hence all
radicands in the following argument are positive. The folded identity is
\[
 P=\delta Z^2-M^2GC_y.
\]
Also \(Q=\Delta^2U^2S>c^2\) and \(F_2>S\), estimates used below.
The three required coprimality facts follow directly from the definitions:
\begin{align*}
 &\gcd(S,F_2)=1,
 &&\gcd(S,\delta C_y)=1,\\
 &\delta=\Delta F_2-A^2,\quad F_2\equiv1\pmod A,
 &&\gcd(F_2,\delta)=1,\\
 &C_y\equiv-\Delta c^2-1=-D_m\pmod{F_2},
 \quad F_2\equiv1\pmod{D_m},
 &&\gcd(F_2,C_y)=1.
\end{align*}
The first line uses \(S\mid Q\); the second uses \(A\mid U\); the
last uses \(D_m\mid S\). No coprimality between \(E\) and \(D_m\)
is asserted.

At a zero, \(\delta z^2=GC_y\). Prime valuations and the preceding
coprimalities force \(G\) to be a positive square. Since its positive
factors \(S,F_2\) are coprime, write
\[
 S=\sigma^2,\qquad F_2=f^2,\qquad G=g^2,
 \qquad g=\sigma f>0.
\]
Then \(y=|z|/g\) is a positive rational solution of
\begin{equation}
 QV^2-(Q-1)y^2=1.\label{eq:compactadjacent}
\end{equation}
It is integral. Indeed, if its reduced denominator is \(m\), then
\(m^2\mid Q-1\). For \(v=|V|>1\), one has
\(v<y<Qv/(Q-1)\), and the transformation
\[
 v'=(2Q-1)v-2(Q-1)y,\qquad
 y'=(2Q-1)y-2Qv
\]
preserves the norm and the reduced denominator, has integer \(v'\), and
satisfies \(0<|v'|<v\). Reflect signs and descend to \(|v|=1\), where
\(|y|=1\). Thus \(m=1\), and \(z=\pm gy\) is integral.

Put \(I=U\sigma\). The recovered strong pair satisfies
\[
 f^2-\Delta I^2=1,\qquad Q=(\Delta I)^2,\qquad c^2\mid I.
\]
The retained all-index folding theorem now applies, giving
\(c=\psi_A(p)\) for a positive odd index \(p\). Its hypotheses are the
strong norm, \(0<c\le I\), the integral auxiliary norm
\eqref{eq:compactadjacent}, and \(V\equiv-c\pmod f\). To see the
folding directly, write \(f=\chi_A(m)\), \(I=\psi_A(m)\). Auxiliary
Pell classification yields an odd index \(j\). The standard odd-quotient
identity, reduced modulo \(f\), gives \(\psi_A(j)\equiv\pm c\pmod f\).
Reducing \(j\) modulo \(2m\) and reflecting gives \(0\le p\le m\).
Both \(\psi_A(p)\) and \(c\) are at most \(I<f/2\), so the congruence
cannot wrap and \(c=\psi_A(p)\); parity makes \(p\) odd.

Consequently
\[
 D_m=\chi_A(p)^2=d^2,\qquad E=(\sigma/d)^2.
\]
An integer that is a rational square is an integer square. Hence
\(E=\tau^2\), \(\sigma=\tau d\), and \(g=\tau df\).
This order is essential: the first square is recovered \emph{after}
main folding, not used as an assumption in that folding.

\subsection{Transport, calibration, and the converse}

The first-Pell classification now gives
\(k=2\psi_{2XY^2+1}(n_0)\). Its first index is impossible, because it
would put \(c=2Y+1\) strictly between the first two main ordinates.
Thus
\[
 n_0\ge2,\qquad k\ge8XY^2+4,\qquad c>kY\ge k>8q^7,
 \qquad c>X,
\]
so \(K>0\). The new strong-size estimate is
\[
 \frac{a^2}{g^2}=\Delta U^2\frac a{1+a}
   >\frac{\Delta U^2}{2}>36.
\]
Together with \(a>2c\), it gives \(|V|>3g\) for nonzero \(b\),
\(\operatorname{sgn}V=\operatorname{sgn}b\), and
\(y<3|V|/2\). The case \(b=0\) would make
\((c^2-1)/(Q-1)=y^2-c^2\) an integer strictly between zero and one.

The parent's third-main-index exclusion uses only the unchanged outer
data: if \(p=3\), then \(n_0=2\), and \(0<\eta<k\) forces
\(2Y-X\in\{1,2\}\), contrary to \(q\mid X,Y\). Thus
\(c>4\Delta+3\). In the auxiliary norm, odd indices 1 and 3 are then
excluded. The two signs at index 3 would require
\[
 F_2(b-4\Delta)=c-(4\Delta+3)
 \quad\hbox{or}\quad
 F_2(b+4\Delta)=c+(4\Delta+3),
\]
whose right sides lie strictly between zero and \(2c<F_2\).
Therefore \(|V|>8Q^2\), \(|b|>4F_2>4g\), and
\(gy<3bV/8<L/2\).

For proof notation let
\(\kappa=K/H\), \(\rho=e/(XY)\), \(t=B/h\), and \(W=1+hL\).
The exact carrier is
\[
 z=n+\kappa-Wt-\rho.
\]
Since \(e=b-cT+k-1<b\), one has \(\rho<|b|/(XY)\).
If \(B\le-1\), the carrier and the preceding estimates would give
\(n<L/2-L+|b|/(XY)<0\). Hence \(B\ge0\), which forces \(C>0\).
The inherited mask bound is now licensed: \(3q<R<q^4-2\).
The normalized rank theorem and the signed auxiliary-index theorem
apply to the recovered strong pair, \(c>2p,2R\), and the two residues
\(V\equiv-c\pmod f\), \(V\equiv-R\pmod c\); they yield \(p=R\).

The rest of the retained carrier proof is unchanged. Integrality of \(z\)
gives \(XY\mid hHe\); its divisor \(q^4\), coprime to \(hH\), gives
\(q^4\mid e\). The first recurrence gives
\(k\equiv2n_0\pmod{4XY^2}\). Since \(n_0<R\), the error
\(2n_0-R-1\) is a multiple of \(q^4\) of absolute value below \(q^4\).
Thus \(2n_0=R+1\), and only then the full recurrence proves
\(XY\mid e\). The retained ratio bound and whole-product projection
yield \(X=2^R\), followed by \(h\mid B\). The fixed-minus parity,
half-binomial rounding, and fresh unscaled native reconstruction therefore
have exactly their parent hypotheses and yield a finite ZFC refutation.
The complete inherited proof closure and its coefficient identities are
included alongside \texttt{compact/PROOF.md}.

Conversely, a refutation supplies the parent's canonical positive outer
witnesses, with \(R\equiv3\pmod4\), \(c=\psi_A(R)\),
\(d=\chi_A(R)\), \(E=\tau^2\), and integral
\(\kappa>0,\rho>0,t\ge0\). Choose one positive strong Pell pair
\((f,I)\) with
\[
 Ac^2\tau d\mid I.
\]
It exists by a finite period of the determinant-one Pell matrix modulo
that integer. Set \(i=I/(Ac^2\tau d)>0\). This realizes exactly
\(F_2=f^2\), \(Q=(\Delta I)^2\), and \(G=(\tau df)^2\).
The retained positive auxiliary reconstruction, which needs only
\(c\mid I\), supplies positive \(T,V,y\) for this \emph{same} strong
pair. Finally choose
\[
 n=gy+(1+hL)t+\rho-\kappa.
\]
Here \(\kappa<Xc<f\le g\), so \(n>0\). The carrier equals \(gy\),
proving \(P=0\). Both existential directions are thus established at
the stated inherited calibration boundary.

\subsection{Exact degree and finite verification}

The evaluator has 46 multiplications, 18 additions, and 14 subtractions;
squares and scalar products count. The decisive exact degrees are
\[
 \deg(Q,G,V,M,Z)=(102,134,96,13,117).
\]
Thus \(-QM^2GV^2\) uniquely supplies degree 454, whereas
\(QZ^2\) has degree 336. Writing \(q_0=\beta J\) and
\(k_0=\eta+\zeta\), the entire highest homogeneous part is
\[
 [P]_{454}=-16q_0^{284}w^{42}s^{80}i^8k_0^{38}T^2.
\]
The coefficient of \(J^{284}w^{42}s^{80}i^8\eta^{38}T^2\) is
\(-16\beta^{284}\ne0\) at the actual fixed tuple.
The independent checker covers all 78 gates by 31 exact polynomial cuts,
reconstructs every coefficient byte, checks all 39 terms of this leading
part, and rejects nine mutations that preserve both gate counts and
liveness. Normal and optimized Python runs pass. These finite checks
authenticate the object; the unbounded proof and inherited calibration
remain the mathematical basis of the consistency equivalence.
